\documentclass[11pt]{article}
\usepackage[a4paper,margin=25mm]{geometry}
\usepackage{amsmath,amssymb,amsthm,mathtools,bm}
\usepackage{booktabs,longtable,array}
\usepackage{enumitem}
\usepackage{xcolor}
\usepackage{microtype}
\usepackage[hidelinks]{hyperref}

\setlength{\emergencystretch}{3em}
\newtheorem{theorem}{Theorem}[section]
\newtheorem{proposition}[theorem]{Proposition}
\newtheorem{lemma}[theorem]{Lemma}
\newtheorem{corollary}[theorem]{Corollary}
\theoremstyle{definition}
\newtheorem{definition}[theorem]{Definition}
\newtheorem{gate}[theorem]{Fail-closed gate}
\theoremstyle{remark}
\newtheorem{remark}[theorem]{Remark}

\newcommand{\ii}{\mathrm i}
\newcommand{\dd}{\mathrm d}
\newcommand{\CC}{\mathbb C}
\newcommand{\ZZ}{\mathbb Z}
\newcommand{\CP}{\mathbb {CP}}
\newcommand{\cO}{\mathcal O}
\newcommand{\cH}{\mathcal H}
\newcommand{\Tr}{\operatorname{Tr}}
\newcommand{\Sym}{\operatorname{Sym}}
\newcommand{\Span}{\operatorname{Span}}
\newcommand{\ket}[1]{\lvert #1\rangle}
\newcommand{\bra}[1]{\langle #1\rvert}
\newcommand{\braket}[2]{\langle #1\vert #2\rangle}
\newcommand{\status}[1]{\textcolor{blue!60!black}{\textsf{#1}}}
\newcolumntype{L}[1]{>{\raggedright\arraybackslash}p{#1}}

\title{AP-E3 UV Follow-up:\\
An Exact Two-Constituent Gauge/Unit-Cell Rule, a Physical \(k=+2\) Orientation,\\
and an Anomaly-Consistent Mixed-WZW Completion Candidate}
\author{Route F research note}
\date{15 July 2026}

\begin{document}
\maketitle

\begin{abstract}
This note closes two sharply delimited AP-E3 questions and exposes the gates
that remain.  First, two independent compact constraint gauge fields impose
\(G_r=N_r-1=0\), \(r=1,2\), on every indivisible unit cell.  The physical
cell Hilbert space is therefore exactly \(\CC^2\otimes\CC^2\); singleton,
odd-total, and charge-transfer sectors are absent kinematically rather than
merely heavy.  Ferromagnetic exchange selects the triplet, while the sign of
the electron Pauli magnetic moment selects its anti-aligned coherent ground
line.  If \(A_+=-\ii s_+^\dagger\dd s_+\) is the AP-E1 unit-flux connection,
then
\[
 \ii\hbar\braket{\Omega_-}{\dd\Omega_-}=+2\hbar A_+.
\]
Thus the signed worldline level is physically \(k=+2\).  The ground line is
the square of the universal quotient line,
\(Q^{\otimes2}\simeq\cO(2)\), and directly has three holomorphic sections.

Second, we construct an anomaly-consistent intermediate UV candidate with
gauge group \(SU(2)_c\times U(1)_g\), two vectorlike Dirac flavours of
abelian charge one, and a charge-one scalar doublet.  All dynamical local and
Witten gauge anomalies cancel.  Gauging \(U(1)_g\) removes the problematic
off-diagonal Pauli--G\"ursey generators, and generalized-symmetry matching
fixes the mixed-WZW coefficient to
\(n=n_cX_q=2\).  A positive unit baryon reduces the five-dimensional term to
\(k=nB=+2\).  The same colour Gauss law forces a local baryon to contain two
quarks and two conjugate-Higgs dressings, producing a degree-two triplet.

The deterministic audit passes all stated algebraic and numerical checks.
Nevertheless, neither branch proves the Route-E family portal.  The mixed-WZW
model still needs nonperturbative vacuum selection, soliton stability,
compact-\(U(1)\)/bordism tests, and an all-scale completion beyond its
abelian Landau pole.  Hence the exact constrained-cell theorem and signed
orientation are closed within their declared microscopic model, while full
AP-E3 UV physics and physics promotion remain false.
\end{abstract}

\tableofcontents

\section{Logical target and result ledger}

The earlier finite-\(U\) Mott--Hund model proves a level magnitude only after
declaring two orbitals and an interacting occupancy plateau.  Finite energy
penalties never remove states from a UV Hilbert space.  The first task here is
therefore not to increase a Mott gap, but to replace the energetic statement
by an exact physical-state condition.  The second task is to select the sign
of the Berry action from a microscopic coupling rather than relabeling an
already derived negative level.  The third task is an independent
four-dimensional topological mechanism with a complete anomaly ledger.

The conclusions must be read at three distinct levels.
\begin{enumerate}[label=(L\arabic*)]
\item \textbf{Exact constrained-cell theorem.}  Once the two compact Gauss
  laws and complete-cell boundary condition are part of the microscopic
  definition, exactly two constituents and the signed \(k=+2\) ground line
  are proved.
\item \textbf{Anomaly-consistent intermediate UV candidate.}  The charged
  two-colour theory cancels every dynamical gauge anomaly and matches the
  required two-group coefficient, but its strong phase and all-scale UV
  completion are not proved.
\item \textbf{Route-E interpretation.}  No calculation below identifies the
  triplet with three chiral \(Spin(10)\) families or derives the portal map.
  Physics promotion therefore remains forbidden.
\end{enumerate}

\section{Exactly two as a gauge/physical-state rule}

\subsection{Why the two Gauss laws are the minimal strong rule}

Four superficially similar restrictions are inequivalent.
\begin{longtable}{@{}L{0.24\textwidth}L{0.21\textwidth}L{0.46\textwidth}@{}}
\toprule
Rule & Physical sector & Decisive boundary\\
\midrule
\endfirsthead
\toprule
Rule & Physical sector & Decisive boundary\\
\midrule
\endhead
\(G_r=N_r-1=0\), \(r=1,2\) & exactly \((N_1,N_2)=(1,1)\) & removes odd and
charge-transfer sectors kinematically; selected here\\
\(G=N_1+N_2-2=0\) & total number two & still contains \((2,0)\) and
\((0,2)\); insufficient without another selector\\
\((-1)^{N_1+N_2}=+1\) & every even sector & contains vacuum, \(N=4,6,\ldots\)
and does not select the pair\\
finite \(\Lambda_G\sum_r(N_r-1)^2\) & all sectors, separated in energy & is
an EFT approximation and cannot be called an exact rule\\
\bottomrule
\end{longtable}

The selected cell contains two co-located but internally labelled local
spin-\(1/2\) moments.  The label may be an orbital, layer, or microscopic
dimer label, but a permitted spatial boundary must terminate whole cells.
This is a model input: the constraint derives the consequences of the cell,
not why nature chose this cell architecture.

\subsection{Constraint action, projector, and large gauge invariance}

For cell \(c\), internal label \(r=1,2\), and spinor index \(A=1,2\), let
\(f_{crA}\) be \emph{bosonic} Schwinger partons with
\begin{equation}
 [f_{crA},f_{c'r'B}^\dagger]
 =\delta_{cc'}\delta_{rr'}\delta_{AB},\qquad
 N_{cr}=f_{crA}^\dagger f_{crA},\qquad
 G_{cr}=N_{cr}-1.
 \label{eq:gauss-law}
\end{equation}
The repeated spinor index is summed.  A worldline action that implements the
constraint is
\begin{equation}
 L_c=\ii\hbar\sum_{r,A}f_{crA}^\dagger\dot f_{crA}
 +\hbar\sum_r a_{0,cr}(N_{cr}-1)-H_c.
 \label{eq:constraint-action}
\end{equation}
Under
\begin{equation}
 f_{cr}\longmapsto e^{\ii\alpha_{cr}}f_{cr},\qquad
 a_{0,cr}\longmapsto a_{0,cr}+\dot\alpha_{cr},
 \label{eq:constraint-gauge}
\end{equation}
the Lagrangian changes only by
\(-\hbar\sum_r\dot\alpha_{cr}\).  On a time circle,
\(\Delta\alpha_{cr}=2\pi m_{cr}\) gives
\begin{equation}
 \exp\left(\frac{\ii\Delta S}{\hbar}\right)
 =\exp\left(-2\pi\ii\sum_r m_{cr}\right)=1.
 \label{eq:large-gauge}
\end{equation}
Thus the integer background charge one has no \(0+1\)-dimensional large-gauge
obstruction.  A half-integer background would fail this test.

The exact physical projector is
\begin{equation}
 \Pi_c^{\rm phys}=\prod_{r=1}^2\int_0^{2\pi}
 \frac{\dd\lambda_{cr}}{2\pi}
 \exp\!\left[\ii\lambda_{cr}(N_{cr}-1)\right].
 \label{eq:exact-projector}
\end{equation}
It follows immediately that
\begin{equation}
 \cH_c^{\rm phys}=\Pi_c^{\rm phys}\cH_c
 =\ker G_{c1}\cap\ker G_{c2}
 \simeq\CC^2\otimes\CC^2,
 \qquad \dim\cH_c^{\rm phys}=4.
 \label{eq:physical-cell}
\end{equation}
All of
\((0,0),(1,0),(0,1),(2,0),(0,2)\) lie outside this Hilbert space.  In
particular, exactly two is an operator identity on physical states,
not the location of a minimum of a potential.

The gauge-invariant spin operators
\begin{equation}
 \bm S_{cr}=\frac12f_{crA}^\dagger\bm\sigma_{AB}f_{crB}
 \label{eq:spin-operator}
\end{equation}
commute with every \(G_{cr'}\).  Elitzur's theorem forbids a parton order
parameter \(\langle f\rangle\), not the physical spin, its spectral projector,
or a gauge-invariant Berry holonomy.

\subsection{Relation to the old finite-Mott candidate}

On the old \((N_1,N_2)=(1,1)\) sector, the map
\begin{equation}
 f_{1A}^\dagger f_{2B}^\dagger\ket0_f
 \longmapsto
 a_{1A}^\dagger a_{2B}^\dagger\ket0_a
 \label{eq:sector-isomorphism}
\end{equation}
intertwines both spin algebras.  Consequently the Hund spectrum, triplet
projector, coherent-state Berry connection, and Veronese map agree exactly.
The difference is entirely UV-logical: the old Mott construction retains
singleton and doublon states at finite energy, whereas
Eq.~\eqref{eq:gauss-law} removes them from the physical Hilbert space.
Recent two-band calculations illustrate why a finite-\(U\) quasiparticle or
local-moment regime must not be confused with an exact projection
\cite{MazitovKatanin2024}.

\subsection{Complete-cell boundary theorem and its limitation}

Let \(P_{0,c}(\bm n)=
\ket{\Omega_{-,c}(\bm n)}\!\bra{\Omega_{-,c}(\bm n)}\) be the unique local
ground projector constructed below and \(Q_{0,c}=1-P_{0,c}\).  If an allowed
finite region \(\Lambda\) consists of complete
cells and an intercell parent interaction satisfies
\begin{equation}
 V_{\rm safe}\ge0,\qquad
 V_{\rm safe}\bigotimes_{c\in\Lambda}\ket{\Omega_{-,c}}=0,
 \qquad [V_{\rm safe},G_{cr}]=0,
 \label{eq:safe-parent}
\end{equation}
then the local gap theorem below implies
\begin{equation}
 H_\Lambda-E_0|\Lambda|
 \ge h\sum_{c\in\Lambda}Q_{0,c}+V_{\rm safe}.
 \label{eq:parent-gap}
\end{equation}
The product ground state is unique and \(\Delta_\Lambda\ge h\).  Thus this
safe parent cone has neither a literal cut-cell singleton nor an accidental
zero-energy boundary mode.

\begin{gate}[Arbitrary interacting boundaries]
The Gauss laws exclude literal odd parton number, but do not forbid every
emergent projective edge representation.  A sufficiently different
intercell interaction may enter a Haldane-like phase and create boundary
spin-\(1/2\) modes without violating \(N_1=N_2=1\).  Outside
Eq.~\eqref{eq:safe-parent}, a gapped homotopy or a direct boundary-spectrum
proof is mandatory.  Lieb--Schultz--Mattis-type constraints are a useful
warning, not a proof for this particular model \cite{YaoEtAl2024}.
\end{gate}

\section{A physical orientation and the signed level \texorpdfstring{\(k=+2\)}{k=+2}}

\subsection{The ferromagnetic sign from exchange}

For two localized electronic orbitals \(\varphi_1,\varphi_2\), the Coulomb
exchange integral can be written
\begin{equation}
 K_{12}=\int\frac{\dd^3q}{(2\pi)^3}
 \frac{4\pi e^2}{\epsilon q^2}
 \left|\widetilde{\varphi_1^*\varphi_2}(\bm q)\right|^2\ge0.
 \label{eq:exchange-integral}
\end{equation}
The direct Coulomb contribution is common to singlet and triplet, whereas
exchange gives
\begin{equation}
 E_{\rm singlet}-E_{\rm triplet}=2K_{12},\qquad J_H=2K_{12}>0.
 \label{eq:hund-sign}
\end{equation}
This supplies a first-principles sign for the effective interaction
\(-J_H\bm S_1\cdot\bm S_2\), subject to the declared localized-orbital
approximation.

\subsection{The Pauli-Zeeman sign is the orientation selector}

For an electron with charge \(-e\) and \(g_{\rm eff}>0\),
\begin{equation}
 \bm\mu=-g_{\rm eff}\mu_B\bm S,
 \qquad H_Z=-\bm\mu\cdot\bm B
 =+g_{\rm eff}\mu_BB_{\rm mag}\,\bm n\cdot\bm S,
 \qquad B_{\rm mag}=|\bm B|,\qquad
 \bm n=\frac{\bm B}{B_{\rm mag}}.
 \label{eq:pauli-sign}
\end{equation}
Absorbing the units of \(\bm S\) into
\(h=g_{\rm eff}\mu_BB_{\rm mag}>0\), the cell
Hamiltonian is
\begin{equation}
 H_c(\bm n)=-J_H\bm S_{c1}\cdot\bm S_{c2}
 +h\bm n\cdot(\bm S_{c1}+\bm S_{c2}),
 \qquad J_H,h>0.
 \label{eq:oriented-hamiltonian}
\end{equation}
The plus sign is fixed by the negative electronic charge.  It is not a manual
reversal of the effective Berry term.

Using
\(\bm S_1\cdot\bm S_2=\tfrac12(\bm S_{\rm tot}^2-\tfrac32)\), the exact
spectrum is
\begin{align}
 E_{1,m}&=-\frac{J_H}{4}+hm,\qquad m=-1,0,+1,\label{eq:triplet-spectrum}\\
 E_{0,0}&=\frac{3J_H}{4}.\label{eq:singlet-spectrum}
\end{align}
The unique ground state is the anti-aligned triplet \(m=-1\), with
\begin{equation}
 E_0=-\frac{J_H}{4}-h,\qquad
 \Delta_{\rm loc}=\min\{h,2h,J_H+h\}=h.
 \label{eq:local-gap}
\end{equation}
Equivalently, as an operator on \(\cH_c^{\rm phys}\),
\begin{equation}
 H_c-E_0\ge h\bigl(1-P_{0,c}(\bm n)\bigr).
 \label{eq:local-gap-operator}
\end{equation}

\subsection{Anti-aligned eigenline as the universal quotient}

Using the standard adiabatic eigenline construction
\cite{Berry1984APE3UV}, let
\begin{equation}
 s_+(\theta,\phi)=
 \begin{pmatrix}\cos(\theta/2)\\ e^{\ii\phi}\sin(\theta/2)\end{pmatrix},
 \qquad
 A_+=-\ii s_+^\dagger\dd s_+
 =\frac{1-\cos\theta}{2}\dd\phi.
 \label{eq:aligned-spinor}
\end{equation}
An eigenvector of \(\bm n\cdot\bm\sigma\) with eigenvalue \(-1\) is
\begin{equation}
 s_-=\ii\sigma_y s_+^*
 =\begin{pmatrix}e^{-\ii\phi}\sin(\theta/2)\\-\cos(\theta/2)\end{pmatrix},
 \qquad -\ii s_-^\dagger\dd s_-=-A_+
 \label{eq:anti-spinor}
\end{equation}
up to an exact patch-gauge term.

There is a useful geometric subtlety.  Let
\begin{equation}
 0\longrightarrow S=\cO(-1)
 \longrightarrow\underline{\CC}^{2}
 \longrightarrow Q\longrightarrow0
 \label{eq:universal-sequence}
\end{equation}
be the universal sequence.  The Hermitian orthogonal complement \(S^\perp\)
is smoothly and complex-linearly isomorphic to the quotient
\(Q\simeq\cO(1)\).  The normalized embedded frame \(s_-=\ii\sigma_y s_+^*\)
is anti-holomorphic in the affine coordinate, but that frame does not define
the holomorphic structure.  Transporting the quotient holomorphic structure
gives \(S^\perp\simeq Q\).

On the north--south overlap,
\begin{equation}
 s_{+,S}=e^{-\ii\phi}s_{+,N},\qquad
 s_{-,S}=e^{+\ii\phi}s_{-,N},
 \label{eq:anti-transition}
\end{equation}
and algebraic quotient frames obey \(q_S=-wq_N\).  Therefore the pair line is
\begin{equation}
 L_-=Q^{\otimes2}\simeq\cO(2).
 \label{eq:positive-pair-line}
\end{equation}

\subsection{Berry action, Chern number, and three states}

The anti-aligned ground state is
\begin{equation}
 \ket{\Omega_-(\bm n)}=\ket{s_-}_1\otimes\ket{s_-}_2.
 \label{eq:anti-pair}
\end{equation}
Tensor-product differentiation gives
\begin{equation}
 \braket{\Omega_-}{\dd\Omega_-}=2s_-^\dagger\dd s_-,
 \qquad
 \boxed{\ii\hbar\braket{\Omega_-}{\dd\Omega_-}=+2\hbar A_+.}
 \label{eq:positive-action}
\end{equation}
Relative to \(S_k=\hbar k\int A_+\), this is the signed result
\begin{equation}
 \boxed{k=+2.}
 \label{eq:kplus2}
\end{equation}
The action connection \(\mathcal B=\ii\braket{\Omega_-}{\dd\Omega_-}\)
has
\begin{equation}
 \dd\mathcal B=2\dd A_+,\qquad
 \frac1{2\pi}\int_{\CP^1}\dd\mathcal B=2.
 \label{eq:positive-chern}
\end{equation}
Consequently positive-polarization quantization directly gives
\begin{equation}
 H^0(\CP^1,L_-)=H^0(\CP^1,\cO(2))
 =\Span\{z_0^2,z_0z_1,z_1^2\},
 \qquad \dim=3.
 \label{eq:positive-sections}
\end{equation}
This corrects a possible overstatement: the normalized orthogonal eigenvector
is anti-holomorphic in \(w\), but its quotient holomorphic line is still
\(\cO(1)\), so the pair is the ordinary positive line \(\cO(2)\).

\begin{gate}[Portal degree]
If Route E supplies an effective orientation map
\(\widehat{\bm B}_{\rm eff}:\CP^1_E\to S^2\), then pullback gives
\begin{equation}
 k_E=2\deg(\widehat{\bm B}_{\rm eff}).
 \label{eq:portal-degree}
\end{equation}
The present calculation proves the coefficient two on the microscopic
orientation sphere, not \(\deg=+1\) for an unknown Route-E portal.
\end{gate}

\section{An anomaly-consistent mixed-WZW intermediate UV candidate}

\subsection{Field content}

Consider the renormalizable gauge theory
\begin{equation}
 G_{\rm gauge}=SU(2)_c\times U(1)_g
 \label{eq:gauge-group}
\end{equation}
with \(N_f=2\) vectorlike Dirac flavours and a scalar doublet.  In an all-left
Weyl notation the relevant representations are
\begin{longtable}{@{}L{0.20\textwidth}L{0.20\textwidth}L{0.22\textwidth}L{0.27\textwidth}@{}}
\toprule
Field & \(SU(2)_c\) & \(U(1)_g\) & global representation\\
\midrule
\endfirsthead
\toprule
Field & \(SU(2)_c\) & \(U(1)_g\) & global representation\\
\midrule
\endhead
\(q_L^{ai}\) & \(\bm2\) & \(+1\) & \(\bm2\) of \(SU(2)_L\)\\
\(q_R^{c\,ai}\) & \(\bm2\) & \(-1\) & \(\overline{\bm2}\) of \(SU(2)_R\)\\
\(\phi_I\) & \(\bm1\) & \(+1\) & \(\bm2\) of \(SU(2)_\phi\)\\
\bottomrule
\end{longtable}
Here \(a=1,2\), \(i=1,2\), and \(I=1,2\).  A potential
\(V_\phi=\lambda(\phi^\dagger\phi-v^2)^2\) produces
\begin{equation}
 \{\phi^\dagger\phi=v^2\}/U(1)_g\simeq S^3/U(1)\simeq\CP^1.
 \label{eq:scalar-cp1}
\end{equation}
The QCD-like condensate is required to break
\(SU(2)_L\times SU(2)_R\to SU(2)_V\), with pion field
\(g\in SU(2)\).

\subsection{Why pure two-colour QCD is a negative control}

Without \(U(1)_g\), pseudoreality of the colour fundamental enlarges the
flavour group to \(SU(2N_f)\) and the standard condensate gives
\begin{equation}
 SU(2N_f)\longrightarrow Sp(2N_f).
 \label{eq:pg-breaking}
\end{equation}
At \(N_f=2\), \(SU(4)/Sp(4)\simeq S^5\) and
\(\pi_3(S^5)=0\).  There is then no stable Skyrmion number to insert into
\(k=n_cB\).  This is why simply substituting \(n_c=2\) into an ordinary
QCD formula is invalid.  Recent two-colour effective studies explicitly
retain the Pauli--G\"ursey structure and diquark directions
\cite{SuenagaEtAl2024APE3UV}.

Gauging \(U(1)_g\) changes this symmetry statement exactly.  In the
Pauli--G\"ursey basis
\(Q=(q_L,\epsilon_cq_R^c)\), the gauge generator is
\begin{equation}
 T_g=\operatorname{diag}(\bm1_{N_f},-\bm1_{N_f}).
 \label{eq:pg-generator}
\end{equation}
Its centralizer is
\begin{equation}
 C_{SU(2N_f)}(T_g)=S\bigl(U(N_f)_L\times U(N_f)_R\bigr).
 \label{eq:pg-centralizer}
\end{equation}
Every off-diagonal Pauli--G\"ursey generator changes \(U(1)_g\) charge by
two and is not a global current.  This removes the symmetry obstruction, but
does not by itself prove which condensate strong dynamics selects.

\subsection{Complete dynamical gauge-anomaly ledger}

The perturbative cubic \(SU(2)_c^3\) anomaly vanishes because the fundamental
is pseudoreal.  The remaining local gauge-anomaly coefficients are
\begin{align}
 \mathcal A_{SU(2)_c^2U(1)_g}
 &=N_fT(\bm2)(1-1)=0,\label{eq:su2su2u1}\\
 \mathcal A_{SU(2)_cU(1)_g^2}
 &=2N_f\bigl(1^2+(-1)^2\bigr)\Tr T^a=0,
 \label{eq:su2u1u1}\\
 \mathcal A_{U(1)_g^3}
 &=2N_f\bigl(1^3+(-1)^3\bigr)=0,\label{eq:u1cube}\\
 \mathcal A_{{\rm grav}^2U(1)_g}
 &=2N_f(1-1)=0.\label{eq:gravu1}
\end{align}
There are \(2N_f=4\) left-Weyl colour doublets, so the mod-two Witten anomaly
also vanishes.  Scalars do not contribute to fermionic anomalies.  Thus the
dynamical gauge theory passes both perturbative and \(SU(2)\) global anomaly
tests.

Background flavour fields reveal the generalized-symmetry datum rather than
a gauge inconsistency.  Let \(b_g\) be the dynamical \(U(1)_g\) connection,
\(h_g=\dd b_g\), and let \(F_{L,R}^{\rm fl}\) be the two background flavour
curvatures.  Here \(n_c=2\) and the quark charge is \(X_q=1\).  In this
normalization, the mixed part of the six-form anomaly polynomial contains
\begin{equation}
 I_6\supset\frac{n_cX_q}{2}\frac{h_g}{2\pi}
 \left[
 \Tr\!\left[\left(\frac{F_L^{\rm fl}}{2\pi}\right)^2\right]
 -\Tr\!\left[\left(\frac{F_R^{\rm fl}}{2\pi}\right)^2\right]
 \right],
 \label{eq:anomaly-polynomial}
\end{equation}
so the Postnikov classes are
\begin{equation}
 \widehat\kappa_L=-\widehat\kappa_R=n_cX_q=2.
 \label{eq:postnikov-two}
\end{equation}
Equation~\eqref{eq:anomaly-polynomial} is written in a Bardeen scheme that
preserves the two dynamical gauge symmetries.  Its background-flavour
variation is compensated by the transformation of the magnetic one-form
background; this is the two-group datum, not a residual gauge anomaly.
Each background \(SU(2)_{L,R}\) also sees two colour copies, an even number,
so its Witten anomaly is trivial.  The scalar \(SU(2)_\phi\) acts only on
bosons.  Bordism refinements remain necessary for the discrete axial and
faithful-group quotients \cite{Saito2024APE3UV,LeeOhmoriTachikawa2021}.

\subsection{Global definition and coefficient of the mixed WZW term}

Whenever the physical four-manifold, the sigma-model field, and the relevant
bundles extend over a five-manifold \(Y_5\), choose integrally normalized forms
\begin{equation}
 \int_{S^3}\omega_3(g)=1,
 \qquad
 \omega_2(s)=\frac{F_s}{2\pi},
 \qquad \int_{S^2}\omega_2=1,
 \label{eq:integral-forms}
\end{equation}
where the trace convention in
\(\omega_3=(24\pi^2)^{-1}\Tr(g^{-1}\dd g)^3\) is fixed by the first equation.
The action is
\begin{equation}
 S_{\rm mix}=2\pi\hbar n\int_{Y_5}\omega_3(g)\wedge\omega_2(s),
 \qquad n\in\ZZ.
 \label{eq:mixed-action}
\end{equation}
Two such extensions differ by a closed five-manifold.  The difference in
\(S/2\pi\hbar\) is \(n\) times an integer, so the exponentiated action is
extension-independent.  Non-extendible sectors require a differential-character
or equivalent \v{C}ech/bordism definition; constructing that refinement remains
an explicit global-definition gate here.  The UV two-group matching fixes,
rather than merely permits,
\begin{equation}
 \boxed{n=n_cX_q=2.}
 \label{eq:wzw-level-two}
\end{equation}
This coefficient is quantized and tree-level exact in the construction of
Davighi and Lohitsiri \cite{DavighiLohitsiri2024APE3UV}.

The local sign can also be followed through the tree-level matching.  At
energies below the Higgs scale, neglecting the \(b_g\) kinetic term at leading
derivative order gives
\begin{align}
 \mathcal L[b_g]
 &=v^2b_{g,\mu}b_g^\mu
 +b_{g,\mu}\bigl(j_\phi^\mu+X_qj_q^\mu\bigr),\notag\\
 b_g^\mu&=-\frac{j_\phi^\mu+X_qj_q^\mu}{2v^2},\qquad
 \mathcal L_{\rm eff}
 =-\frac{(j_\phi+X_qj_q)^2}{4v^2}.
 \label{eq:bg-eom}
\end{align}
Confinement identifies \(j_q=n_cj_B\).  Locally define
\(A_\phi=j_\phi/(2v^2)=+\ii s^\dagger\dd s=-A_R\), where
\(A_R=-\ii s^\dagger\dd s\) is our Route-E/AP-E1 convention and
\(F_s=\dd A_R\).  The cross term is therefore
\begin{equation}
 S_{\rm cross}=-n_cX_q\hbar\int *j_B\wedge A_\phi
 =+n_cX_q\hbar\int *j_B\wedge A_R.
 \label{eq:local-cross}
\end{equation}
Thus the microscopic current coupling fixes the local sign; the
five-dimensional expression supplies its global, quantized completion.

\subsection{Worldline reduction and physical orientation}

For a baryon/Skyrmion with
\begin{equation}
 B=\int_{M_3}\omega_3(g)\in\ZZ,
 \label{eq:baryon-number}
\end{equation}
take \(Y_5=M_3\times D_2\), with \(\partial D_2=\gamma\).  Product
orientation and Stokes' theorem give
\begin{align}
 S_{\rm mix}
 &=2\pi\hbar n
 \left(\int_{M_3}\omega_3\right)
 \left(\int_{D_2}\frac{F_s}{2\pi}\right)\notag\\
 &=\hbar nB\oint_\gamma A_R.
 \label{eq:worldline-reduction}
\end{align}
Therefore
\begin{equation}
 \boxed{k=nB=2B.}
 \label{eq:k-nb}
\end{equation}
The positively oriented baryon sector \(B=+1\) has \(k=+2\); its CPT
antibaryon has \(B=-1\) and \(k=-2\).  This orientation statement is
physical but sector-dependent, not a claim of a net CPT asymmetry.

\subsection{Colour Gauss law, Higgs dressing, and the even-level rule}

One colour fundamental is not a local gauge singlet.  The minimal local
baryon contains two quarks and carries \(U(1)_g\) charge \(+2\).  A local
gauge-invariant interpolating operator must therefore contain two conjugate
Higgs fields:
\begin{equation}
 \mathcal B^{IJ}_{ij}\sim
 \epsilon_{ab}(q_i^{aT}Cq_j^b)
 \phi^{\dagger (I}\phi^{\dagger J)}.
 \label{eq:dressed-baryon}
\end{equation}
The scalar factor is
\begin{equation}
 \Sym^2\overline{\bm2}=\bm3,
 \qquad
 \cO(1)^{\otimes2}=\cO(2).
 \label{eq:dressing-triplet}
\end{equation}
The antisymmetric same-point scalar singlet vanishes.  Thus the first local
colour-singlet baryon is simultaneously an exactly-two scalar dressing, a
degree-two line, and an \(SU(2)_\phi\) triplet.

There is also a faithful-group explanation of evenness.  The complete triple
\[
 (-\bm1_c,e^{\ii\pi}_g,-\bm1_\phi)
\]
acts trivially on both quarks and scalars.  Equivalently, the global centre
\(-\bm1_\phi\) agrees on all fields with the gauge transformation
\((-\bm1_c,e^{\ii\pi}_g)\).  Hence the faithful global action on local
gauge-invariant scalar operators is
\begin{equation}
 SU(2)_\phi/\ZZ_2\simeq SO(3).
 \label{eq:faithful-so3}
\end{equation}
An \(SO(3)\)-equivariant \(\cO(k)\) line requires \((-1)^k=1\).  The local
physical sector therefore has even \(k\), with minimal nonzero value two.
Nonlocal line defects are an explicit exception and may carry projective
representations.

\section{What remains before this is a full UV completion}

\subsection{Strong phase and scale ordering}

The charged theory removes the exact Pauli--G\"ursey symmetry, but a desired
vacuum still has to be selected dynamically.  The required inequalities are
schematically
\begin{equation}
 m_{\rm PG}^2>0,\qquad
 m_{\rm PG}\gg E_{\rm worldline},\qquad
 \Lambda_c\gtrsim m_b\sim e_gv.
 \label{eq:phase-gates}
\end{equation}
The first two prevent the baryon from unwinding through an approximately
\(S^5\) diquark direction.  The last ordering keeps confinement in a regime
where \(U(1)_g\) still distinguishes the two Pauli--G\"ursey blocks.  If
\(v\gg\Lambda_c\), the heavy vector leaves only suppressed interactions and
the unwanted symmetry may become approximately accurate.  These questions
need lattice or another controlled nonperturbative calculation.

\subsection{Running and the all-scale obstruction}

With two Dirac flavours, two colours, and two charge-one complex scalars, the
one-loop coefficients are
\begin{equation}
 b_{U(1)}=\frac43(2N_f)+\frac13(2)=6,
 \qquad
 b_{SU(2)_c}=\frac{11}{3}(2)-\frac43T(\bm2)N_f=6.
 \label{eq:beta-coefficients}
\end{equation}
Colour is asymptotically free, but the abelian coupling has a Landau pole.
The theory is therefore a valid anomaly-consistent intermediate UV
description above the sigma model, not an all-scale completion.

A proposed nonabelian embedding
\(SU(3)\to[SU(2)_c\times U(1)]/\ZZ_2\) has a sharp representation-theoretic
cost: every \(SU(2)_c\) singlet descending from an \(SU(3)\) representation
has even abelian charge.  A light colour-singlet charge-one \(\phi_I\) must
therefore arrive with coloured partners, whose decoupling must preserve both
light scalars and \(SU(2)_\phi\).  No such complete decoupling proof is
provided here.

One creative research branch replaces \(U(1)_g\) by an asymptotically free
\(SU(2)_g\), broken by an adjoint to its Cartan.  Bifundamental quarks and
scalar \(SU(2)_g\) doublets could reproduce the intermediate spectrum, but
mass splitting makes the chiral symmetry emergent and
't Hooft--Polyakov monopoles make the magnetic one-form symmetry only
approximate.  It is a future branch, not a claimed solution.

\subsection{Soliton and quantization gates}

Even a quantized \(k=2\) term does not prove a three-state low-energy particle.
The following remain independent calculations:
\begin{enumerate}
\item existence and metastability of a \(B=1\) soliton in the full field
  space, including its Hessian and decay channels;
\item Finkelstein--Rubinstein statistics and the collective-coordinate
  moment of inertia;
\item separation of the \(\cO(2)\) lowest sector from the higher rotor tower;
\item compact-\(U(1)\) monopole breaking of the magnetic one-form symmetry;
\item discrete axial and spin-bordism consistency;
\item a Route-E portal whose degree is \(+1\) and whose couplings remain below
  every retained gap.
\end{enumerate}

\section{Deterministic regression and numerical evidence}

The companion verifier
\texttt{route\_f/code/verify\_ap\_e3\_exact\_two\_mixed\_wzw.py}
performs the following independent checks.
\begin{enumerate}
\item It enumerates the bosonic four-mode Fock basis truncated at
  \(0\le N_r\le2\) and evaluates the exact Fourier projector
  \eqref{eq:exact-projector}.  The separate projector has rank four; within
  this declared audit truncation, the single total-number projector has rank
  ten and retains both charge-transfer sectors.  The exact constrained rank
  four is truncation-independent.
\item It diagonalizes \eqref{eq:oriented-hamiltonian} at \(J_H=1,h=0.2\).
  The spectrum is
  \begin{equation}
  (-0.45,-0.25,-0.05,0.75),
  \label{eq:numerical-spectrum}
  \end{equation}
  and seeded random directions agree with the anti-aligned tensor-square
  projector to floating-point precision.
\item Exact directional derivatives verify
  \(\ii\braket{\Omega_-}{\dd\Omega_-}=2A_+\) and doubled quantum metric.
  Midpoint curvature quadrature converges monotonically to \(c_1=+2\), while
  the overlap transition has winding two.
\item The script evaluates every local anomaly sum, the mod-two colour and
  flavour counts, \(\widehat\kappa_L=-\widehat\kappa_R=2\), extension
  independence of the exponentiated five-dimensional action, and
  \(k(B=+1)=+2\).
\item Negative controls include one total Gauss law, parity-only projection,
  a half-integer background charge, pure two-colour QCD, and the abelian
  Landau pole.
\end{enumerate}
These tests validate formulas and conventions; they cannot replace the
nonperturbative phase and soliton calculations listed above.

\section{Claim ledger and ordered continuation}

\begingroup
\small
\begin{longtable}{@{}L{0.40\textwidth}L{0.19\textwidth}L{0.32\textwidth}@{}}
\toprule
Claim & Status & Boundary\\
\midrule
\endfirsthead
\toprule
Claim & Status & Boundary\\
\midrule
\endhead
Separate compact Gauss laws impose \(N_1=N_2=1\) & \status{proved} & within
the declared indivisible cell\\
Literal odd and charge-transfer sectors are absent & \status{proved} & not a
finite-energy statement\\
Safe-parent complete-cell boundaries have gap at least \(h\) &
\status{proved} & arbitrary intercell phases excluded\\
Pauli orientation gives signed \(k=+2\) & \status{proved} & on the microscopic
orientation sphere\\
Positive ground line is \(Q^2=\cO(2)\) with three sections &
\status{proved} & family interpretation not implied\\
Dynamical gauge anomalies of the charged two-colour model cancel &
\status{proved} & discrete/bordism refinements remain\\
Mixed-WZW coefficient and \(B=1\) reduction give \(k=+2\) &
\status{proved conditionally} & extendible sectors; assumes the desired gapped
strong phase and unit soliton\\
Global differential-cohomology definition on non-extendible sectors &
\status{open} & extension independence alone is insufficient\\
All-scale mixed-WZW UV completion & \status{open} & abelian Landau pole and
nonabelian decoupling problem\\
Route-E degree-one chiral-family portal & \status{open} & no microscopic map
or anomaly-complete family operator\\
\bottomrule
\end{longtable}
\endgroup

The next steps should be executed in this order:
\begin{enumerate}
\item lattice or another nonperturbative audit of the charged two-colour
  phase, measuring meson/diquark condensates, \(m_{\rm PG}\), and the
  \(B=1\) lifetime over \((e_g,v/\Lambda_c)\);
\item APS/bordism and compact-monopole audit of the faithful global group;
\item collective-coordinate Hessian, statistics, and rotor/LLL gap;
\item explicit degree-one Route-E portal below all gaps;
\item use the now-complete AP-E4 mathematical spectrum as an independent
  control, then derive its still-open physical fermion origin through either
  moduli-space SQM or an anomaly-complete product compactification.
\end{enumerate}

The fail-closed status is therefore
\begin{equation}
 \boxed{\begin{gathered}
 \texttt{exactly\_two\_constrained\_cell\_proved=true},\\
 \texttt{physical\_pauli\_orientation\_gives\_k\_plus\_2=true},\\
 \texttt{dynamical\_gauge\_anomalies\_cancel=true},\\
 \texttt{mixed\_wzw\_intermediate\_uv\_candidate=true},\\
 \texttt{mixed\_wzw\_full\_uv\_closed=false},\\
 \texttt{route\_e\_portal\_closed=false},\\
 \texttt{physics\_promotion\_allowed=false}.
 \end{gathered}}
 \label{eq:final-status}
\end{equation}

\bibliographystyle{unsrt}
\bibliography{route_f/tex/ap_e3_exact_two_mixed_wzw_uv}

\end{document}
